\documentclass[10pt,a4paper]{amsart}
\usepackage[margin=19mm]{geometry}
\usepackage{amsmath,amssymb,mathtools}
\usepackage[scr=rsfs,cal=euler]{mathalfa}
\usepackage{xfrac}
\usepackage[unicode,hidelinks,bookmarks=false]{hyperref}
\newcommand{\ie}{i.e.\ }
\newcommand{\cf}{cf.\ }
\newcommand{\resp}{respectively\ }
\newcommand{\Subsection}{\subsection}
\providecommand{\nobreakdash}{\nobreak\hbox{-}}
\setlength{\emergencystretch}{2em}
\multlinegap0pt
\usepackage{bm}
\usepackage{bbm}
\usepackage{dsfont}
\usepackage{xspace}
\usepackage{cancel}
\usepackage{mathtools}
\usepackage{amsthm}
\makeatletter
\@ifundefined{theorem}{\newtheorem{theorem}{Theorem}}{}
\@ifundefined{lemma}{\newtheorem{lemma}[theorem]{Lemma}}{}
\@ifundefined{proposition}{\newtheorem{proposition}[theorem]{Proposition}}{}
\makeatother
\def \R {\mathbb{R}}
\def \So {\mathscr{S}(\Rm)}
\newcommand{\Rm}{\mathbb R^m}
\newcommand{\da}{\rangle}
\newcommand{\p}{\partial}
\newcommand{\sa}{\langle}
\newcommand{\tr}{\operatorname{tr}}
\newcommand{\vf}{\varphi}
\title[Schr\"odinger semigroups and the H\"ormander hypoellipticity]{Schr\"odinger semigroups and the H\"ormander hypoellipticity condition}
\author{Nicola Garofalo, Alessandra Lunardi}
\date{Source: arXiv:2406.04441v1 (CC BY 4.0)}
\begin{document}
\maketitle

\noindent\emph{Source and licence.} Schr\"odinger semigroups and the H\"ormander hypoellipticity condition. Version arXiv:2406.04441v1 (CC BY 4.0), distributed under the Creative Commons Attribution 4.0 International licence, as recorded in the arXiv licence metadata for this exact version and shown on the abstract page https://arxiv.org/abs/2406.04441v1. Full rights notice: licenses/arxiv-2406.04441-CC-BY-4.0.txt.\par\smallskip

\noindent\textit{Editorial scope.} Counted unit: the Proposition whose source label is \texttt{P:erwin} in the article (source0.tex lines 541--555, source label \texttt{P:erwin}), delivered together with the article's own complete original proof of it (source0.tex lines 558--616, one \texttt{proof} environment of 59 lines). No mathematics is written, repaired or re-derived: statement and proof are the article's own text line for line. Required context retained uncounted: A (lines 332--338); Ds (lines 340--342); f (lines 346--354); L:drift (lines 450--464); drift (lines 452--454); Abis (lines 456--462); P:ftgaussian (lines 514--527); gengaussi2 (lines 516--520); imgaussian (lines 522--526). Every definition, display and label of those lines is retained unchanged. Omitted from this package and not redistributed: 1--331, 339--339, 343--345, 355--449, 455--455, 463--513, 521--521, 527--540, 556--557, 617--1321. Nothing inside a retained excerpt is omitted or altered: every retained body is delivered exactly as the article's lines are written, so no omission receipt of any kind accompanies this package. Citations: the retained excerpts carry no citation command at all, so no bibliography entry of the article is transcribed and no reference is redistributed. Reference adaptation: none was needed and none was made; every reference inside the delivered unit resolves inside it, so no unresolved reference or citation is produced. Printed slips kept literal: no printed word, symbol, hypothesis or number of the counted statement or of its proof was altered. Layout and class: the article's own class is \texttt{amsart} with packages including amssymb, amsmath, epsfig, graphics, mathrsfs, enumerate, verbatim, hyperref, fancyhdr, amsfonts, amsthm, graphicx, dsfont, color; the wrapper is the lane's pinned \texttt{amsart} editorial wrapper at 10pt on A4, which restates the theorem-like environments the retained text needs and adds the article's own macro definitions that the retained text needs (\texttt{\textbackslash R}, \texttt{\textbackslash Rm}, \texttt{\textbackslash So}, \texttt{\textbackslash da}, \texttt{\textbackslash p}, \texttt{\textbackslash sa}, \texttt{\textbackslash tr}, \texttt{\textbackslash vf}), with extra wrapper packages (bm, bbm, dsfont, xspace, cancel, mathtools). The delivered numbering is the unit's own. Assets: no retained span contains a figure environment, a graphics-inclusion command, a PDF-page inclusion, a table or a TikZ picture; no image file is copied into this package, the package declares no asset dependency and no image file is redistributed with it. Licence: version arXiv:2406.04441v1 of this article is distributed under the Creative Commons Attribution 4.0 International licence, as recorded in the arXiv licence metadata for this exact version and shown on the abstract page https://arxiv.org/abs/2406.04441v1. A full rights notice is delivered at licenses/arxiv-2406.04441-CC-BY-4.0.txt. Characters: every retained excerpt is pure ASCII, so the delivered engine, XeLaTeX with its default Latin Modern text font, reports no missing character.
\begin{equation}\label{A}
\begin{cases}
\p_t f- i \operatorname{tr}(Q \nabla^2 f) -  \langle B x,\nabla f\rangle  = 0,
\\
f(x,0) = \vf(x),\ \ \ \ \ \vf\in \mathscr S(\Rm),
\end{cases}
\end{equation} 

\begin{equation}\label{Ds}
Q(t) = \int_0^t e^{sB}Q e^{sB^\star} ds > 0.
\end{equation}

\begin{equation}\label{f}
\mathcal T(t) \vf(x)  = \begin{cases}
(4\pi
)^{-\frac{m}{2}} \frac{e^{-\frac{i \pi m}4}}{\sqrt{\det Q(t)}}   \int_{\Rm} e^{i \frac{\sa Q(t)^{-1}(y-e^{tB}x),y-e^{tB}x\da}{4}} \vf(y) dy,\ \ \ \ t>0,
\\
\\
\vf(x),\ \ \ \ \ \ \ \ \ \ \  t = 0,
\end{cases}
\end{equation}

\begin{lemma}\label{L:drift}
Let $Q, B\in \mathbb M_{m\times m}(\R)$, and suppose that $v$ and $f$ are connected by the relation
\begin{equation}\label{drift}
v(x,t) = f(e^{- t B} x,t).
\end{equation}
Then, $f$ is a solution of the Cauchy problem \eqref{A} if and only if $v$ solves the problem
\begin{equation}\label{Abis}
\begin{cases}
\p_t v - i \operatorname{tr}(Q'(t) \nabla^2 v) = 0,
\\
v(x,0) = \vf(x),
\end{cases}
\end{equation} 
where $Q(t)$ is the matrix defined in \eqref{Ds}.
\end{lemma}

\begin{proposition}\label{P:ftgaussian}
Let $A\in G\ell(\mathbb C,m)$ be such that $A^\star = A$ and $\Re A \ge 0$. Then 
\begin{equation}\label{gengaussi2}
\mathscr F\left(\frac{1}{\sqrt{\operatorname{det} A}}(4\pi
)^{-\frac{m}{2}} e^{- \frac{\sa A^{-1}\cdot,\cdot\da}{4}}\right)(\xi) =
e^{- 4 \pi^2  \sa A\xi,\xi\da},
\end{equation}
where $\sqrt{\operatorname{det} A}$ is the unique analytic branch such that $\sqrt{\operatorname{det} A}>0$ when $A$ is real. In particular, when $A = i A_0$, with $A_0$ real, symmetric and $>0$, then \eqref{gengaussi2} gives
\begin{equation}\label{imgaussian}
\mathscr F\left(\frac{e^{-\frac{i m\pi}4}}{\sqrt{\operatorname{det} A_0}}(4\pi
)^{-\frac{m}{2}} e^{\frac{i \sa A_0^{-1}\cdot,\cdot\da}{4}}\right)(\xi) =
e^{- 4 \pi^2 i  \sa A_0\xi,\xi\da}.
\end{equation}
\end{proposition}

\begin{proposition}\label{P:erwin}
In $\Rm\times \Rm\times (0,\infty)$ consider the kernel
\begin{equation}\label{Skolmo}
\mathcal S(x,y,t) = 
(4\pi
)^{-\frac{m}{2}} \frac{e^{-\frac{i m\pi}4}}{\sqrt{\det Q(t)}}   e^{i \frac{\sa Q(t)^{-1}(y-x),y-x\da}{4}}\ \ \ \ t>0.
\end{equation}
Given $\vf\in \mathscr S(\Rm)$, the function
\begin{equation}\label{reallyfinal}
f(x,t) = 
\int_{\Rm} \mathcal S(e^{t B} x,y,t) \vf(y) dy,\ \ \ \ \ t>0,
\end{equation}
and $f(x,0) = \vf(x)$, 
is the unique solution of the Cauchy problem \eqref{A} such that $f(\cdot,t)\in \mathscr S(\Rm)$ for every $t\ge 0$.
\end{proposition}

\begin{proof}
Given $\vf\in \So$, let $f$ solve \eqref{A}. According to Lemma \ref{L:drift}, the function $v$ defined by \eqref{drift} solves the problem \eqref{Abis}. To obtain a representation of $v$, we take a partial Fourier transform with respect to $x\in \Rm$,
\[
\hat v(\xi,t) = \int_{\Rm} e^{-2\pi i\sa \xi,x\da} v(x,t) dx.
\]
The problem \eqref{Abis} is thus changed into
\begin{equation}\label{AbisFT}
\begin{cases}
\p_t \hat v + 4\pi^2 i \sa Q'(t) \xi,\xi\da \hat v = \p_t \hat v + 4\pi^2 i \frac{d}{dt} \sa Q(t) \xi,\xi\da \hat v = 0,
\\
\hat v(\xi,0) = \hat \vf(\xi).
\end{cases}
\end{equation} 
We obtain from \eqref{AbisFT}
\begin{equation}\label{hatv}
\hat v(\xi,t) = \hat \vf(\xi) e^{-4\pi^2 i \sa Q(t) \xi,\xi\da},
\end{equation}
where in the second equality we have used the definition (but not the positivity, yet) of the covariance matrix $Q(t)$ in \eqref{Ds}. To proceed, we next want to show that, for every $t>0$, the function $\xi\to e^{-4\pi^2 i  \sa Q(t) \xi,\xi\da}$ is a Fourier transform. 
For this we use \eqref{imgaussian}  in Proposition \ref{P:ftgaussian} with $A_0 = Q(t)>0$, obtaining
\begin{equation}\label{imHor}
\mathscr F\left(\frac{e^{-\frac{i m\pi}4}}{\sqrt{\det{Q(t)}}}(4\pi
)^{-\frac{m}{2}} e^{i \frac{\sa Q(t)^{-1}\cdot,\cdot\da}{4}}\right)(\xi) =
e^{- 4 \pi^2 i \sa Q(t)\xi,\xi\da}.
\end{equation}
From \eqref{hatv} and \eqref{imHor} we obtain for $t>0$ the following representation for the solution of the Cauchy problem \eqref{Abis} 
\begin{equation}\label{tildeTt}
v(x,t)  = (4\pi
)^{-\frac{m}{2}} \frac{e^{-\frac{i m\pi}4}}{\sqrt{\det Q(t)}} \int_{\Rm}  e^{i \frac{\sa Q(t)^{-1}(x-y),x-y\da}{4}} \vf(y) dy = \int_{\Rm} \mathcal S(x,y,t) \vf(y) dy,
\end{equation}
where in the second equality we have used \eqref{Skolmo}.
Combining \eqref{drift} with \eqref{tildeTt}, we finally obtain \eqref{f}, or equivalently \eqref{reallyfinal}. 

To complete the proof, and also for later purposes, we prove a representation of the spatial Fourier transform of $f(x,t) = \mathcal T(t) \vf(x)$, were $\mathcal T(t)$ is defined by \eqref{f}. With this in mind, we recall the following well-known property. If $g\in L^1(\Rm)$ and $A\in G\ell(\R,m)$, then one has
\begin{equation}\label{ftA}
\hat g(A^\star \xi) = |\operatorname{det} A|^{-1} \widehat{g \circ A^{-1}}(\xi).
\end{equation} 
Applying \eqref{ftA} with $A = e^{-tB}$, and keeping in mind that $\det A = e^{-t \tr B}$, we obtain
\[
\mathscr F(v(e^{t B} \cdot,t))(\xi) = e^{-t \tr B} \hat v(e^{-tB^\star}\xi).
\]
From this identity and from  \eqref{f}, we obtain the following basic representation formula
\begin{equation}\label{eccola}
\widehat{f(\cdot,t)}(\xi) = e^{-t \tr B} \hat \vf(e^{-tB^\star}\xi) \  e^{-4\pi^2 i  \sa Q(t) e^{-tB^\star}\xi,e^{-tB^\star}\xi\da}.
\end{equation}
Since it is clear that 
\[
\xi\ \longrightarrow\ \hat \vf(e^{-tB^\star}\xi) \in \mathscr S(\Rm),
\]
and one easily verifies that the $C^\infty$ function 
\[
\xi\ \longrightarrow\  e^{-4\pi^2 i  \sa Q(t) e^{-tB^\star}\xi,e^{-tB^\star}\xi\da}
\]
is a multiplier for $\mathscr S(\Rm)$, we conclude from \eqref{eccola} that $\widehat{f(\cdot,t)}\in \mathscr S(\Rm)$ for every $t\ge 0$, and therefore $f(\cdot,t) \in \mathscr S(\Rm)$. From these considerations, and from the (almost) obvious consequence $\widehat{f(\cdot,t)}\to \hat \vf$ of \eqref{eccola}, we conclude that 
\[
f(\cdot,t)\ \underset{t\to 0^+}{\longrightarrow}\ \vf.
\]
To complete the proof of Proposition \ref{P:erwin}, we are left with showing that \eqref{f} solves the PDE in \eqref{A}. Since we know that the function $f$ in \eqref{f} satisfies $f(\cdot,t)\in \So$, the fact that it is a solution to the PDE in \eqref{A} follows from tracing back the steps that lead to \eqref{f}. Finally, its uniqueness is again a direct consequence of \eqref{eccola}.

\end{proof}

\end{document}
